% !TeX root = ../教材题库_全部条目_转写工作稿.tex
% Source: XJ-B1-C04-SU-P150-Y1441
% 第4章 / 小结与复习 小结与复习 / 教材 p.150 / PDF 第157页
% OCR未校对：公式与文字均可能有误；未完成数学TeX转写；禁止直接发给学生
% 原题及相邻原解截图（核对用）：../题目原图/XJ-B1-C04-SU-P150-Y1441-01.webp
% 本文件只在转写完成后作为题干使用；OCR 条目尚不是可靠题干。
\begin{bankitem}{第4章 小结与复习 第7题}{第4章 / 小结与复习 小结与复习 / 教材 p.150 / PDF 第157页}{XJ-B1-C04-SU-P150-Y1441}{OCR未校对：公式与文字均可能有误；未完成数学TeX转写；禁止直接发给学生}
\begin{ocrdraft}
[教材p.150 / PDF第157页]
\par
7．用函数的观点环视大千世界，我们看到许多事物发展过程中的某些具有
\par
标志性的数量可以看成时间的函数，能用适当的函数构建数学模型来描述．收
\par
集一些社会生活中普遍使用的函数模型（指数函数、对数函数、幂函数以及分段
\par
函数等）的实例，了解函数模型的广泛应用．
\end{ocrdraft}
\end{bankitem}
